\documentclass[letterpaper,11pt]{article}

\usepackage[empty]{fullpage}
\usepackage[hidelinks]{hyperref}
\usepackage{fontawesome5}
\usepackage{enumitem}
\usepackage{parskip}

% Adjust margins
\addtolength{\oddsidemargin}{-0.4in}
\addtolength{\evensidemargin}{-0.4in}
\addtolength{\textwidth}{0.8in}
\addtolength{\topmargin}{-0.5in}
\addtolength{\textheight}{1.0in}

\ifdefined\pdfgentounicode
  \input{glyphtounicode}
  \pdfgentounicode=1
\fi

\begin{document}
\begin{center}
    {\Huge \scshape Juan Becerra} \\ \vspace{4pt}
    \textbf{\large Software Engineer -- Gameplay, Animation \& Systems} \\ \vspace{4pt}
    \small Seattle, WA ~ $|$ ~
    \raisebox{-0.1\height}\faPhone\ (208) 550-9191 ~ $|$ ~
    \href{mailto:becerrajuan007@gmail.com}{\raisebox{-0.2\height}\faEnvelope\ \underline{becerrajuan007@gmail.com}} ~ $|$ ~
    \href{https://hoowan.dev}{\raisebox{-0.2\height}\faGlobe\ \underline{hoowan.dev}} \\ \vspace{2pt}
    \href{https://linkedin.com/in/hoowan/}{\raisebox{-0.2\height}\faLinkedin\ \underline{linkedin.com/in/hoowan}} ~ $|$ ~
    \href{https://github.com/hoowan-dev}{\raisebox{-0.2\height}\faGithub\ \underline{github.com/hoowan-dev}}
\end{center}

\vspace{4pt}
\hrule
\vspace{14pt}

\noindent November 2, 2026

\vspace{8pt}
\noindent \textbf{Nintendo Software Technology Hiring Team} \\
Nintendo Software Technology Corporation \\
Redmond, WA

\vspace{8pt}
\noindent \textbf{Re: CONTRACT - Gameplay Software Engineer}

\vspace{8pt}
\noindent Dear Nintendo Software Technology Hiring Team,

I am writing to express my strong interest in the \textbf{Contract Gameplay Software Engineer} position at
Nintendo Software Technology. Over the past five years at Electronic Arts, as a Gameplay \& Animation Software
Engineer at \textbf{Maxis} on \emph{Project Rene} (\emph{The Sims}) and now as a Tools \& Tech-Art Software
Engineer at \textbf{EA Lifestyle}, I have built player-facing gameplay systems in C++ and worked side by side
with designers and animators to make them feel right. Nintendo's games are the reason I care so much about
responsive controls and polish, and I would love to contribute to that standard at NST.

At \textbf{EA Maxis Studios}, my individual contributions centered on how characters move and respond to the
player. I implemented core locomotion, player-facing character controls, state machines, and runtime animation
systems in Unreal Engine 5 and C++, including motion matching, distance matching, and custom IK solvers. I also
architected a multi-threaded, asynchronous path-planning and navigation API that delivered deterministic
character routing and real-time obstacle avoidance within a frame budget of under 1.5ms. Working in a large,
multi-disciplinary codebase meant reading, debugging, and extending other engineers' code every day, and
profiling and optimizing memory and frame budgets to keep gameplay free of hitches. Throughout, I partnered
closely with gameplay designers, technical animators, and QA, building in-editor tuning tools and visual
debuggers so we could prototype, iterate on, and polish features quickly together.

At \textbf{EA Lifestyle}, I engineered a high-throughput Unreal Engine C++ plugin that captures runtime character
trajectory, pose, and velocity data for motion-matching databases, and built a Maya/Python stylization tool
adopted across multiple studios that cut redundant keyframing work by over 70\%. Earlier, as a Senior Games
Test Engineer at \textbf{Lionbridge}, I led regression testing on live-service AAA titles, which gave me a deep
respect for the quality bar a game has to meet before it reaches players.

Outside of studio work, I keep building games end to end: \emph{DoomSlop}, a 3D first-person shooter with
responsive movement and camera controls, hitscan weapons, boss encounters, and a gameplay HUD;
\emph{BroncoDrome}, a team-built vehicle combat game in Unreal Engine 4 using C++ and Blueprints; \emph{Chasm},
an action RPG in Unity and C\#; and a custom C++ engine with an entity-component system for game objects.

Based in Seattle, I would be glad to work hybrid in Redmond and bring clear, maintainable C++ and C\# gameplay
code to NST's team. Thank you for your time and consideration. I look forward to speaking with you.

\vspace{10pt}
\noindent Sincerely,

\vspace{14pt}
\noindent \textbf{Juan Becerra}

\end{document}
